\paragraph{Per-ball win-probability and expected-runs models.}
The methodological backbone of in-play cricket modeling is a dynamic-programming
value function over remaining balls and wickets, $V(b,w) = r(b,w) + p(b,w)\,
V(b{+}1,w{+}1) + (1{-}p(b,w))\,V(b{+}1,w)$, whose ingredients are exactly the two
quantities we target: per-ball expected runs $r(b,w)$ and wicket probability
$p(b,w)$. WASP \citep{brooker2012wasp} popularized this recursion but
deliberately models an average team against an average team, treating player
identity as out of scope. \citet{asif2016inplay} forecast win probability with a
dynamic logistic regression whose coefficients evolve smoothly over an innings,
prizing calibration over identity. \citet{davis2021dynamic} predict
second-innings outcomes ball by ball with engineered per-player metrics, and
\citet{lamsal2024ingame} extend the dynamic-regression line to IPL T20 with
Bayesian priors and Gaussian-process smoothing. \citet{swartz2016simulator,
swartz2020xer} build hierarchical empirical-Bayes per-ball outcome models that
\emph{do} condition on batter and bowler, but for simulation and an
expected-goals-style economy metric rather than a state-versus-identity
decomposition. None of these reports the relative weight of identity against
state, isolated to the death overs and scored properly for both targets.

\paragraph{Pressure indices and contextual performance.}
A second strand operationalizes ``pressure'' as a deterministic function of game
state. \citet{shah2014pressure} introduce a ball-by-ball pressure index from
required rate, wickets, and resources; \citet{bhattacharjee2016quantifying}
refine it into the widely used PI3; \citet{mallawa2024pressure} derive a
target-aware index from a differential equation; and
\citet{thomson2021contextual} define ``clutch'' batting and bowling on a
runs-to-resources scale. A recent higher-order Markov analysis
\citep{arxiv2025markov} confirms that pressure-index variance is largest in the
death overs. Crucially, every member of this family \emph{derives} pressure from
state and then validates it by predicting outcomes; none separates
pressure-as-state from pressure-as-residual-history, which is precisely the
Markov-violation question we test.

\paragraph{Machine learning for cricket, including the death phase.}
A large applied-ML literature predicts match winners or innings totals with
gradient boosting and random forests \citep{raju2024dream11},
typically reporting accuracy or $R^2$ on aggregate quantities rather than
calibrated per-ball probabilities. A few works engage the death phase directly:
\citet{sahoo2024shot} use SHAP to model shot selection in the final over, and
\citet{jamil2023factors} identify death-phase performance factors in
fifty-over cricket. These establish the death overs as a distinct regime and SHAP
as the natural attribution tool, but stop short of a calibrated, properly scored
per-ball model of both expected runs and wicket probability.

\paragraph{The hot hand and the streak-selection correction.}
\citet{gilovich1985hothand} reported no positive autocorrelation in basketball
shooting, founding the view that the hot hand is a cognitive illusion.
\citet{miller2018surprised} proved that the conditional-probability estimator
used in that analysis carries a finite-sequence selection bias: conditioning on a
realized streak inside a finite sequence makes the expected
post-streak success rate strictly below the true rate, with the bias growing in
streak length and shrinking in sequence length. Correcting for it reverses the
canonical result and recovers real hot-hand effects, confirmed across datasets in
\citet{miller2024coldshower}. The tennis-momentum literature
\citep{wang2024momentum} has recently adopted difference-equation and
machine-learning models of within-match momentum. Cricket has its own hot-hand strand, but it does not use the
streak-selection correction: \citet{durbach2007cricket} test the perception of
randomness in cricket scoring sequences with run-based tests that predate
\citet{miller2018surprised}, and \citet{ram2022hothand} report a significant
hot-hand effect across players' \emph{careers} using a self-exciting Hawkes point
process on match-level performances, citing Miller--Sanjurjo only to frame the
broader debate rather than applying their finite-sequence bias correction. Neither
operates at the per-ball level nor de-biases the conditional-probability
estimator. To our knowledge the finite-sequence bias correction itself has not
been applied to cricket at the per-ball level.

\paragraph{What this paper does that prior work does not.}
We combine these strands into one experiment on the death-over regime. Unlike the
DP/dynamic-regression line, we report a state-versus-identity skill decomposition
for both targets with held-out incremental skill, not just attribution. Unlike
the pressure-index family, we isolate a residual history-beyond-state effect and
test the Markov assumption directly. And unlike prior cricket hot-hand work, which
cites Miller--Sanjurjo but does not de-bias the conditional estimator, we apply the
finite-sequence bias correction itself --- with a state-stratified permutation null
--- to per-ball momentum, to our knowledge the first such per-ball application in
the sport.
